% GENERATED FILE. Edit canonical lessons, then run npm run generate:books.
% canonical-source-hash: fnv1a64:5ad260ca2f983e67
% canonical-lessons: MR-C85-masa, MR-C85-bhat, MR-C85-dal, MR-C85-phal, MR-C85-safarchand

\chapter{Fish, Cooked rice, Lentils, Fruit, Apple}
\label{ch:mr-fish-cooked-rice-lentils-fruit-apple}
\hlchaptermodality{85}

\begin{chapteropening}
\textbf{By the end of this chapter:} \emph{I can say a fish, cooked rice, lentils, a fruit and an apple.}

Complete the last lesson of chapter 85: I can say a fish, cooked rice, lentils, a fruit and an apple.
\end{chapteropening}

\section[\texorpdfstring{\mr{मासा} (māsā)}{मासा (māsā)}]{\mr{मासा} (māsā) — a fish}

\label{lesson:MR-C85-masa}



\begin{quote}
Before the new one: say the Marathi for to fall, then the Marathi for meat.
\end{quote}

\subsection*{You'll want to know: \mr{मासा}}
\textbf{\mr{मासा}} — \emph{māsā} — \textquotedblleft{}a fish\textquotedblright{}.

\begin{cousinweb}
\textbf{\mr{मासा}} (\emph{māsā}) is a tadbhava, a word inherited from Sanskrit and reshaped by sound change. It comes from \emph{matsya}, \textquotedblleft{}fish\textquotedblright{} (CDIAL 9758).
\end{cousinweb}

\begin{grammarlens}[title={What you've built}]
One more word for this chapter.
\end{grammarlens}

\subsection*{Your turn}
\begin{itemize}
\raggedright
  \item \emph{Say it:} \emph{māsā}
  \item \emph{Say it:} \emph{māsā}, once more
  \item \emph{Say it:} say \emph{māsā}
  \item \emph{Recall:} say the Marathi for to fall, then the Marathi for meat, then say \emph{māsā} again
\end{itemize}

\begin{tcolorbox}[breakable,colback=teal!4,colframe=teal!35!black,title={Before you move on}]
What does \mr{मासा} mean? (\textquotedblleft{}A fish\textquotedblright{}.) Say it once more.
\end{tcolorbox}
\section[\texorpdfstring{\mr{भात} (bhāt)}{भात (bhāt)}]{\mr{भात} (bhāt) — cooked rice}

\label{lesson:MR-C85-bhat}



\begin{quote}
Before the new one: say the Marathi for meat, then the Marathi for a fish.
\end{quote}

\subsection*{You'll want to know: \mr{भात}}
\textbf{\mr{भात}} — \emph{bhāt} — \textquotedblleft{}cooked rice\textquotedblright{}.

\begin{cousinweb}
\textbf{\mr{भात}} (\emph{bhāt}) is a tadbhava, a word inherited from Sanskrit and reshaped by sound change. It comes from \emph{bhakta}, \textquotedblleft{}meal, food; boiled rice\textquotedblright{} (CDIAL 9331). It is the past participle of \emph{bhaj}, \textquotedblleft{}to share out\textquotedblright{}: a portion.
\end{cousinweb}

\begin{grammarlens}[title={What you've built}]
One more word for this chapter.
\end{grammarlens}

\subsection*{Your turn}
\begin{itemize}
\raggedright
  \item \emph{Say it:} \emph{bhāt}
  \item \emph{Say it:} \emph{bhāt}, once more
  \item \emph{Say it:} say \emph{bhāt}
  \item \emph{Recall:} say the Marathi for meat, then the Marathi for a fish, then say \emph{bhāt} again
\end{itemize}

\begin{tcolorbox}[breakable,colback=teal!4,colframe=teal!35!black,title={Before you move on}]
What does \mr{भात} mean? (\textquotedblleft{}Cooked rice\textquotedblright{}.) Say it once more.
\end{tcolorbox}
\section[\texorpdfstring{\mr{डाळ} (ḍāḷ)}{डाळ (ḍāḷ)}]{\mr{डाळ} (ḍāḷ) — lentils}

\label{lesson:MR-C85-dal}



\begin{quote}
Before the new one: say the Marathi for a fish, then the Marathi for cooked rice.
\end{quote}

\subsection*{You'll want to know: \mr{डाळ}}
\textbf{\mr{डाळ}} — \emph{ḍāḷ} — \textquotedblleft{}lentils\textquotedblright{}.

\begin{cousinweb}
\textbf{\mr{डाळ}} (\emph{ḍāḷ}) is a tadbhava of Sanskrit \emph{dāla}, \textquotedblleft{}a kind of grain, split peas\textquotedblright{} (CDIAL 6309). CDIAL asks whether it is the same word as \emph{*dāla}, \textquotedblleft{}tearing, splitting\textquotedblright{}.
\end{cousinweb}

\begin{grammarlens}[title={What you've built}]
One more word for this chapter.
\end{grammarlens}

\subsection*{Your turn}
\begin{itemize}
\raggedright
  \item \emph{Say it:} \emph{ḍāḷ}
  \item \emph{Say it:} \emph{ḍāḷ}, once more
  \item \emph{Say it:} say \emph{ḍāḷ}
  \item \emph{Recall:} say the Marathi for a fish, then the Marathi for cooked rice, then say \emph{ḍāḷ} again
\end{itemize}

\begin{tcolorbox}[breakable,colback=teal!4,colframe=teal!35!black,title={Before you move on}]
What does \mr{डाळ} mean? (\textquotedblleft{}Lentils\textquotedblright{}.) Say it once more.
\end{tcolorbox}
\section[\texorpdfstring{\mr{फळ} (phaḷ)}{फळ (phaḷ)}]{\mr{फळ} (phaḷ) — a fruit}

\label{lesson:MR-C85-phal}



\begin{quote}
Before the new one: say the Marathi for cooked rice, then the Marathi for lentils.
\end{quote}

\subsection*{You'll want to know: \mr{फळ}}
\textbf{\mr{फळ}} — \emph{phaḷ} — \textquotedblleft{}a fruit\textquotedblright{}.

\begin{cousinweb}
\textbf{\mr{फळ}} (\emph{phaḷ}) is a tadbhava of Sanskrit \emph{phala}, \textquotedblleft{}fruit\textquotedblright{} (CDIAL 9051). CDIAL compares Tamil \emph{paṛam}, \textquotedblleft{}ripe fruit\textquotedblright{}, and Telugu \emph{paṇḍu}, \textquotedblleft{}to ripen\textquotedblright{}.
\end{cousinweb}

\begin{grammarlens}[title={What you've built}]
One more word for this chapter.
\end{grammarlens}

\subsection*{Your turn}
\begin{itemize}
\raggedright
  \item \emph{Say it:} \emph{phaḷ}
  \item \emph{Say it:} \emph{phaḷ}, once more
  \item \emph{Say it:} say \emph{phaḷ}
  \item \emph{Recall:} say the Marathi for cooked rice, then the Marathi for lentils, then say \emph{phaḷ} again
\end{itemize}

\begin{tcolorbox}[breakable,colback=teal!4,colframe=teal!35!black,title={Before you move on}]
What does \mr{फळ} mean? (\textquotedblleft{}A fruit\textquotedblright{}.) Say it once more.
\end{tcolorbox}
\section[\texorpdfstring{\mr{सफरचंद} (safarchand)}{सफरचंद (safarchand)}]{\mr{सफरचंद} (safarchand) — an apple}

\label{lesson:MR-C85-safarchand}



\begin{quote}
Before the new one: say the Marathi for lentils, then the Marathi for a fruit.
\end{quote}

\subsection*{You'll want to know: \mr{सफरचंद}}
\textbf{\mr{सफरचंद}} — \emph{safarchand} — \textquotedblleft{}an apple\textquotedblright{}.

\begin{cousinweb}
The origin of \textbf{\mr{सफरचंद}} (\emph{safarchand}) is not settled. It is often linked with Persian \emph{safarjal}, \textquotedblleft{}quince\textquotedblright{}, from Arabic, but CDIAL has no entry for it and no source here confirms the link.
\end{cousinweb}

\begin{grammarlens}[title={What you've built}]
One more word for this chapter.
\end{grammarlens}

\subsection*{Your turn}
\begin{itemize}
\raggedright
  \item \emph{Say it:} \emph{safarchand}
  \item \emph{Say it:} \emph{safarchand}, once more
  \item \emph{Say it:} say \emph{safarchand}
  \item \emph{Recall:} say the Marathi for lentils, then the Marathi for a fruit, then say \emph{safarchand} again
\end{itemize}

\begin{tcolorbox}[breakable,colback=teal!4,colframe=teal!35!black,title={Before you move on}]
What does \mr{सफरचंद} mean? (\textquotedblleft{}An apple\textquotedblright{}.) Say it once more.
\end{tcolorbox}
